\NSPage{007}%
The estimates must control every Cartesian space-time derivative of this residual.

First we construct the axisymmetric background
\NSi{NS-F-P007-001}{(u_B,p_B)} of Proposition~\ref{prop:5.5}, whose residual
in \eqref{eq:5.41} is the negative divergence of an annular stress, up to a remainder vanishing to every
order at the singular point. Oscillatory velocities supply the leading stress through their
averaged quadratic products (Proposition~\ref{prop:7.5}). Successive corrections improve the decay of
the remaining residual (Proposition~\ref{prop:9.6}). We then sum and localize the fields, preserving
incompressibility and the leading velocity growth, and extend the residual to the smooth
compactly supported force in Theorem~\ref{thm:1.1} (Propositions~\ref{prop:10.1} and~\ref{prop:9.9} and
Lemma~\ref{lem:10.3}). All estimates used below are proved for the viscous equations. The construction
is summarized in Figures~\ref{fig:5} and~\ref{fig:6}.

For any \NSi{NS-F-P007-002}{\nu>0}, a solution at viscosity one gives a solution at viscosity
\NSi{NS-F-P007-003}{\nu} by
\NSFormula{NS-F-P007-004}{display}{none}
\[
 u_\nu(x,t)=\sqrt{\nu}\,u(x/\sqrt{\nu},t),\qquad
 p_\nu(x,t)=\nu p(x/\sqrt{\nu},t),\qquad
 f_\nu(x,t)=\sqrt{\nu}\,f(x/\sqrt{\nu},t).
\]
The singular time is unchanged. The equation and energy scaling are verified in
\eqref{eq:10.22}--\eqref{eq:10.23}. We use \NSi{NS-F-P007-005}{\nu=1} for the rest of this outline.

\subsection{The concentrating leading field.}\label{sec:3.1}
We encode the geometry of Section~\ref{sec:2.1} by the fixed velocity profiles
\NSi{NS-F-P007-006}{E,U} and pressure profile \NSi{NS-F-P007-007}{\Pi} in \eqref{eq:4.3}.
The similarity coordinates \NSi{NS-F-P007-008}{(X,\eta)} in \eqref{eq:3.2} resolve the different
radial and axial contraction rates and express the velocity growth as powers of a concentration
scale. We first specify this representation; the next subsection constructs profiles whose
remaining momentum residual can be supplied by the pulses.

We write
\NSFormula{NS-F-P007-009}{display}{none}
\[
 \tau=1-t,\qquad A=\frac12+h,\qquad D=\frac12-h,
 \qquad 0<h<1/100,
\]
with \NSi{NS-F-P007-010}{h} fixed in the profile construction. In cylindrical coordinates,
\NSFormula{NS-F-P007-011}{display}{none}
\[
 x=(r\cos\theta,r\sin\theta,z),\qquad
 e_r=(\cos\theta,\sin\theta,0),\qquad
 e_\theta=(-\sin\theta,\cos\theta,0),\qquad
 e_z=(0,0,1).
\]
The leading field
\NSi{NS-F-P007-012}{u^{(0)}=u_r^{(0)}e_r+u_\theta^{(0)}e_\theta+u_z^{(0)}e_z}
is axisymmetric: its cylindrical components are independent of \NSi{NS-F-P007-013}{\theta}.
The tangential components are the azimuthal and axial components, tangent to cylinders of
constant \NSi{NS-F-P007-014}{r}.

We introduce the similarity coordinates
\NSFormula{NS-F-P007-015}{display}{3.2}
\begin{equation}
 \tau=q(1-\eta^2),\qquad z=q^D\eta,\qquad X=\frac{r^2}{2q},
 \qquad q>0,\qquad -1<\eta<1.
 \tag{3.2}\label{eq:3.2}
\end{equation}
For \NSi{NS-F-P007-016}{\tau>0}, eliminating \NSi{NS-F-P007-017}{\eta} gives
\NSFormula{NS-F-P007-018}{display}{none}
\[
 q-z^2q^{2h}=\tau,\qquad
 \partial_q(q-z^2q^{2h})=1-2h\eta^2\geq 1-2h>0
 \quad\text{on }|\eta|<1.
\]
Thus \NSi{NS-F-P007-019}{q=q(z,\tau)} is uniquely determined, and
\NSi{NS-F-P007-020}{q\asymp\tau+|z|^{1/D}}, where \NSi{NS-F-P007-021}{\asymp} denotes bounds above
and below by fixed positive factors. On a fixed compact set of similarity coordinates away
from \NSi{NS-F-P007-022}{\eta=\pm1}, we have \NSi{NS-F-P007-023}{q\asymp\tau}.
For bounded \NSi{NS-F-P007-024}{X}, letting \NSi{NS-F-P007-025}{q\downarrow0} approaches the
singular point. The endpoints \NSi{NS-F-P007-026}{\eta=\pm1}, with
\NSi{NS-F-P007-027}{q>0}, describe \NSi{NS-F-P007-028}{t=1} away from that point; see
Lemma~\ref{lem:4.1}.

The azimuthal and axial profiles \NSi{NS-F-P007-029}{E(X,\eta),U(X,\eta)} determine the leading
fields by
\NSFormula{NS-F-P007-030}{display}{none}
\[
 u_\theta^{(0)}=q^{-A}E,\qquad u_z^{(0)}=q^{-A}U,
 \qquad
 ru_r^{(0)}=V_0,\qquad p^{(0)}=q^{-2A}\Pi.
\]
Incompressibility and regularity at the axis determine \NSi{NS-F-P007-031}{V_0}:
\NSFormula{NS-F-P007-032}{display}{none}
\[
 \frac1r\partial_r(ru_r^{(0)})+\partial_z u_z^{(0)}=0,
 \qquad V_0(0,\eta)=0.
\]
\NSPage{008}%
The leading radial pressure balance and the normalization at radial infinity give
\NSFormula{NS-F-P008-001}{display}{none}
\[
 \partial_r p^{(0)}=\frac{(u_\theta^{(0)})^2}{r},
 \qquad
 \Pi(X,\eta)=-\int_X^\infty\frac{E(x,\eta)^2}{2x}\,dx.
\]
The construction ensures that \NSi{NS-F-P008-002}{E/\sqrt{2X}},
\NSi{NS-F-P008-003}{U}, and \NSi{NS-F-P008-004}{V_0/X} are smooth at
\NSi{NS-F-P008-005}{X=0}, giving a smooth Cartesian field at the axis for every
\NSi{NS-F-P008-006}{t<1}. In particular, \NSi{NS-F-P008-007}{u_\theta^{(0)}=0} on the axis,
while \NSi{NS-F-P008-008}{E(X,\eta)>0} for \NSi{NS-F-P008-009}{X>0}.

We choose \NSi{NS-F-P008-010}{X_c>0} and \NSi{NS-F-P008-011}{0<\eta_c<1} so that
\NSi{NS-F-P008-012}{[0,X_c]\times[-\eta_c,\eta_c]} lies in the inner region of
Theorem~\ref{thm:4.6}, and define the core at time \NSi{NS-F-P008-013}{t=1-\tau} by
\NSFormula{NS-F-P008-014}{display}{none}
\[
 \mathcal C_\tau=\{(r\cos\theta,r\sin\theta,z):0\leq X(r,z,\tau)\leq X_c,
 \ |\eta(z,\tau)|\leq\eta_c\}.
\]
Since \NSi{NS-F-P008-015}{q\asymp\tau} there, its radial and axial extents satisfy
\NSFormula{NS-F-P008-016}{display}{none}
\[
 \ell_r\asymp\tau^{1/2},\qquad
 \ell_z\asymp\tau^{1/2-h},\qquad
 \frac{\ell_z}{\ell_r}\asymp\tau^{-h}.
\]
Both extents vanish; their ratio diverges. The constructed profiles give
\NSFormula{NS-F-P008-017}{display}{none}
\[
 \begin{gathered}
  \|u_\theta^{(0)}\|_{L^\infty(\mathcal C_\tau)}\asymp\tau^{-1/2-h},
  \qquad
  \|u_z^{(0)}\|_{L^\infty(\mathcal C_\tau)}\asymp\tau^{-1/2-h},\\
  \|u_r^{(0)}\|_{L^\infty(\mathcal C_\tau)}=O(\tau^{-1/2}),
  \qquad
  \frac{\|u_r^{(0)}\|_{L^\infty(\mathcal C_\tau)}}
       {\|u_\theta^{(0)}\|_{L^\infty(\mathcal C_\tau)}}=O(\tau^h).
 \end{gathered}
\]
The axial lower bound follows from the nonzero axis datum in Proposition~\ref{prop:B.2}.
For any fixed \NSi{NS-F-P008-018}{0<X_*<X_c}, the circle
\NSi{NS-F-P008-019}{z=0,\ r=\sqrt{2X_*\tau}} has
\NSi{NS-F-P008-020}{q=\tau,\eta=0}, and
\NSFormula{NS-F-P008-021}{display}{none}
\[
 u_\theta^{(0)}=E(X_*,0)\tau^{-1/2-h}\longrightarrow+\infty.
\]
The core is a time-dependent region in space; its boundary is determined by fixed similarity
coordinates. Incompressibility preserves the volume of material regions transported by the
velocity. The radial regions are shown in Figure~\ref{fig:3}(a).

The leading field extends beyond the core. Outside a fixed radial interval in
\NSi{NS-F-P008-022}{X}, \NSi{NS-F-P008-023}{U=V_0=0}, while
\NSFormula{NS-F-P008-024}{display}{none}
\[
 E(X,\eta)=c_\infty X^{-1/2-h}\bigl(1+O(X^{-1})\bigr)
 \qquad (X\to\infty),\qquad c_\infty>0,
\]
uniformly in \NSi{NS-F-P008-025}{\eta}. The exact exterior formula \eqref{eq:4.29} gives a purely
azimuthal solution of the radial heat equation for the azimuthal velocity. The final spatial
localization in Proposition~\ref{prop:10.1} gives the completed velocity a fixed compact support.

At viscosity one, the local length scales
\NSi{NS-F-P008-026}{\ell_r=q^{1/2},\ \ell_z=q^D} give
\NSFormula{NS-F-P008-027}{display}{none}
\[
 \operatorname{Re}_\theta=q^{-A}\ell_r=q^{-h}\longrightarrow\infty,
 \qquad
 \operatorname{Re}_r=|u_r^{(0)}|\ell_r=O(1).
\]
The rotation rate exceeds the radial diffusion rate. In the axisymmetric tangential equations,
radial diffusion and radial and axial transport enter the leading balance at the rates
\NSFormula{NS-F-P008-028}{display}{none}
\[
 \frac{|u_r^{(0)}|}{\ell_r}=O(q^{-1}),
 \qquad
 \frac{q^{-A}}{\ell_z}=\frac1{\ell_r^2}=q^{-1}.
\]
Axial diffusion is smaller than radial diffusion by the factor
\NSi{NS-F-P008-029}{\ell_r^2/\ell_z^2=q^{2h}}, which supplies the expansion parameter for the
background corrections.

\NSPage{009}%
\begin{figure}[htbp]
 \centering
 \NSFigurePlaceholder{NS-FIG-3}
 \caption{The concentrating geometry and the reference pulse. (a) A schematic radial--axial
 section at fixed time: the pulse annulus lies between the level surfaces
 \NSi{NS-F-P009-001}{X=X_a} and \NSi{NS-F-P009-002}{X=X_b}. Horizontal and vertical axes use their
 respective similarity scales; both physical length scales shrink. (b) A reference amplitude
 \NSi{NS-F-P009-003}{P(v)}, normalized to one at the midpoint of a pulse interval of length
 \NSi{NS-F-P009-004}{L_s} in \NSi{NS-F-P009-005}{v}, grows and then decays. The shaded regions
 contain the small cutoff tails. The spatial radii and pulse parameters shown are representative;
 the proof chooses them by the estimates in the profile and pulse sections.}
 \label{fig:3}
\end{figure}

\subsection{Construction of the background stress.}\label{sec:3.2}
We now choose the profiles \NSi{NS-F-P009-006}{E,U} so that the pulses can supply their remaining
momentum transport. The resulting stress \NSi{NS-F-P009-007}{T}, defined by the radial integrals
below, must vanish outside the pulse annulus and be a positive combination of the fluxes carried
by the two wave families. The profile construction in Theorem~\ref{thm:4.6}(ii) gives this support;
the admissible stress cone condition characterized in Lemma~\ref{lem:4.5} ensures the positive
representation of Proposition~\ref{prop:7.5}.

We denote by \NSi{NS-F-P009-008}{R_\theta^{(0)},R_z^{(0)}} the tangential momentum residuals of
\NSi{NS-F-P009-009}{(u^{(0)},p^{(0)})}, retaining radial viscosity and omitting axial viscosity.
We construct physical stress components
\NSi{NS-F-P009-010}{T=(T_{r\theta},T_{rz})} satisfying
\NSFormula{NS-F-P009-011}{display}{none}
\[
 R_\theta^{(0)}=-(\partial_r+2/r)T_{r\theta},
 \qquad
 R_z^{(0)}=-(\partial_r+1/r)T_{rz}.
\]
Regularity at the axis fixes their integration constants:
\NSFormula{NS-F-P009-012}{display}{none}
\[
 T_{r\theta}(r)=-\frac1{r^2}\int_0^r s^2R_\theta^{(0)}(s)\,ds,
 \qquad
 T_{rz}(r)=-\frac1r\int_0^r sR_z^{(0)}(s)\,ds.
\]
We require \NSi{NS-F-P009-013}{T} to be nonzero precisely in the annulus
\NSi{NS-F-P009-014}{X_a<X<X_b}, or
\NSi{NS-F-P009-015}{\sqrt{2qX_a}<r<\sqrt{2qX_b}}, with
\NSi{NS-F-P009-016}{0<X_a<X_b}. We make the residual zero in the inner region and in the heat
exterior, and impose
\NSFormula{NS-F-P009-017}{display}{none}
\[
 \int_0^\infty r^2R_\theta^{(0)}\,dr=0,
 \qquad
 \int_0^\infty rR_z^{(0)}\,dr=0.
\]
These identities make the stress zero beyond the exterior radius as well as near the axis. The
prescribed exterior moment identities give these zero integrals (Lemma~\ref{lem:A.8}). Matching
the profiles and the five radial integrals defined in \eqref{eq:4.15} preserves the exterior fields
(Lemma~\ref{lem:4.4}).

\NSPage{010}%
The admissible stress cone condition permits the pulse construction to select two families with
linearly independent covariance vectors \NSi{NS-F-P010-001}{v_1,v_2\in\mathbb R^2} at each point
of the annulus. These vectors depend on the local background and give the azimuthal and axial
momentum fluxes per unit squared amplitude. They represent the required stress as
\NSFormula{NS-F-P010-002}{display}{none}
\[
 T=c_1v_1+c_2v_2,\qquad c_1,c_2>0\quad\text{where }T\neq0.
\]
Thus \NSi{NS-F-P010-003}{T} lies in the interior of the cone generated by
\NSi{NS-F-P010-004}{v_1,v_2}. Proposition~\ref{prop:7.5} constructs this representation from the
admissible stress cone condition on the profiles. Figure~\ref{fig:4} illustrates this positive
representation.

\begin{figure}[htbp]
 \centering
 \NSFigurePlaceholder{NS-FIG-4}
 \caption{Positive representation of the target stress. The shaded cone consists of nonnegative
 linear combinations of the covariance directions \NSi{NS-F-P010-005}{v_1,v_2}. An interior
 target \NSi{NS-F-P010-006}{T=c_1v_1+c_2v_2} has
 \NSi{NS-F-P010-007}{c_1,c_2>0}; the dashed lines show vector addition. The coefficients remain
 positive under sufficiently small perturbations of the two directions.}
 \label{fig:4}
\end{figure}

These conditions couple the inner and exterior profiles through the pressure and radial velocity,
both of which depend on cumulative radial integrals. We choose the pressure datum, solve near the
axis, and match these integrals before imposing the cone condition:
\begin{enumerate}[label=\arabic*.]
 \item \emph{Exterior and axis pressure.} Proposition~\ref{prop:A.4} constructs reference profiles
 with a purely azimuthal tail. Their radial pressure integral \eqref{eq:4.31} fixes
 \NSi{NS-F-P010-008}{\Pi(0,\eta)}. Replacing the tail by the exact heat flow of
 Lemma~\ref{lem:A.6} changes this integral. The corrections in Proposition~\ref{prop:A.7} restore
 it at smaller radii, preserving the axis pressure.

 \item \emph{Inner solution.} With this pressure datum and the specified axis data,
 Proposition~\ref{prop:B.2} constructs \NSi{NS-F-P010-009}{E/\sqrt{2X},U,\Pi} analytic in
 \NSi{NS-F-P010-010}{X,\eta} near the axis and satisfying \eqref{eq:4.13}. Their leading
 tangential residual and stress vanish on \NSi{NS-F-P010-011}{X\leq X_a}.

 \item \emph{Joining and moment matching.} Corollary~\ref{cor:B.10} connects the inner profiles to
 the prescribed outer profiles. The correction in Proposition~\ref{prop:B.8} matches the five
 cumulative radial integrals governing the pressure, radial velocity, and stress. Equality of the
 profiles and these integrals at the joining radius restores the prescribed exterior fields by
 Lemma~\ref{lem:4.4}. The inner identity \eqref{eq:4.13} and the exterior conclusion of
 Lemma~\ref{lem:A.8} give \NSi{NS-F-P010-012}{T=0} for
 \NSi{NS-F-P010-013}{X\leq X_a} and \NSi{NS-F-P010-014}{X\geq X_b}.

 \item \emph{Enforcing the cone condition.} On a compact subinterval of the annulus,
 Proposition~\ref{prop:C.2} introduces a radial oscillation with phase
 \NSi{NS-F-P010-015}{N\log X}, where \NSi{NS-F-P010-016}{N} is a large fixed integer. Its
 amplitude is \NSi{NS-F-P010-017}{O(N^{-1})}; applying
 \NSi{NS-F-P010-018}{X\partial_X} produces an order-one change in the radial derivatives. This
 changes the leading tangential shear while changing profile values and radial moments by only
 \NSi{NS-F-P010-019}{O(N^{-1})}. The shear follows a periodic curve satisfying the admissible
 stress cone condition. A separate localized correction restores all five radial integrals
 exactly. The inner solution and heat exterior are preserved, and the resulting stress admits
 positive covariance weights throughout the annulus (Proposition~\ref{prop:C.3}).
\end{enumerate}

\NSPage{011}%
The leading profiles are now fixed. Before adding pulses, Section~\ref{sec:5} constructs corrections
in successive powers \NSi{NS-F-P011-001}{q^{2nh},\ n=1,2,\ldots}. At each order, the previously
determined coefficients supply the source terms in the coefficient equations. This recursion
accounts for the radial momentum equation, axial viscosity, and the remaining lower-order terms.
The coefficients are summed with smooth cutoffs on their vector potentials and pressures, equal to
one near the singularity and supported in successively smaller neighborhoods. Taking curls after
multiplication preserves incompressibility. The resulting smooth axisymmetric background
\NSi{NS-F-P011-002}{(u_B,p_B)} satisfies
\NSFormula{NS-F-P011-003}{display}{none}
\[
 \begin{aligned}
 \mathscr R(u_B,p_B)
   & =-(\partial_r+2/r)\mathcal T_{\mathrm{phys},\theta}e_\theta \\
   & \phantom{={}}-(\partial_r+1/r)\mathcal T_{\mathrm{phys},z}e_z+\mathcal E_B.
 \end{aligned}
\]
The stress \NSi{NS-F-P011-004}{\mathcal T_{\mathrm{phys}}} is supported in the annulus and has
leading term \NSi{NS-F-P011-005}{T}. A residual is \emph{flat} as
\NSi{NS-F-P011-006}{q\downarrow0} if every Cartesian space-time derivative is
\NSi{NS-F-P011-007}{O(q^N)} for every \NSi{NS-F-P011-008}{N\geq0}, uniformly on each bounded
interval in \NSi{NS-F-P011-009}{X}. The remainder \NSi{NS-F-P011-010}{\mathcal E_B} has this
property by Proposition~\ref{prop:5.5}.

\subsection{Oscillations and momentum transport.}\label{sec:3.3}
We've now arrived at a background \NSi{NS-F-P011-011}{(u_B,p_B)} with the required growth and
leading annular stress \NSi{NS-F-P011-012}{T} that the two pulse families can realize. We next
construct divergence-free pulses (Lemma~\ref{lem:7.7}) whose averaged quadratic flux supplies
\NSi{NS-F-P011-013}{T} to leading order (Propositions~\ref{prop:7.5} and~\ref{prop:9.5}). The exact
increment identity separates this cancellation from the linear pulse evolution and the interactions
left for later corrections.

For a divergence-free velocity increment \NSi{NS-F-P011-014}{w} with pressure increment
\NSi{NS-F-P011-015}{\pi}, the exact residual identity is
\NSFormula{NS-F-P011-016}{display}{none}
\[
 \mathscr R(u_B+w,p_B+\pi)=\mathscr R(u_B,p_B)+\mathcal L_{u_B}(w,\pi)
   +\nabla\cdot(w\otimes w),
\]
where
\NSFormula{NS-F-P011-017}{display}{none}
\[
 \mathcal L_{u_B}(w,\pi):=\partial_t w+(u_B\cdot\nabla)w+(w\cdot\nabla)u_B
   -\Delta w+\nabla\pi.
\]
The linear term describes the evolution of the pulses in the background; the quadratic term
provides the stress that cancels the leading part of
\NSi{NS-F-P011-018}{\mathscr R(u_B,p_B)}.

We write \NSi{NS-F-P011-019}{A_{\mathrm{wave}}} for the characteristic pulse oscillation amplitude
and \NSi{NS-F-P011-020}{\ell_{\mathrm{wave}}} for its wavelength. At fixed nonzero profile values in
the pulse annulus, suppressing logarithmic factors and weights that vanish at the support boundary,
the construction gives
\NSFormula{NS-F-P011-021}{display}{none}
\[
 \begin{gathered}
  A_{\mathrm{wave}}\asymp q^{-1/2-h/2},\qquad
  \ell_{\mathrm{wave}}\asymp q^{1/2+h/2},\\
  \frac{A_{\mathrm{wave}}}{q^{-1/2-h}}\asymp
  \frac{\ell_{\mathrm{wave}}}{q^{1/2}}\asymp q^{h/2}.
 \end{gathered}
\]
The square of the amplitude has the required stress scale:
\NSFormula{NS-F-P011-022}{display}{none}
\[
 A_{\mathrm{wave}}^2\asymp q^{-1-h}\asymp
 \frac{|u_\theta^{(0)}|}{q^{1/2}},
 \qquad
 \frac{A_{\mathrm{wave}}^2}{q^{1/2}}\asymp q^{-3/2-h}.
\]
The last expression is the scale of the averaged radial stress divergence, matching the leading
time derivative, transport, and radial viscosity in the tangential momentum equations.

Freezing the amplitude and phase gradient at a point gives the schematic local form
\NSFormula{NS-F-P011-023}{display}{none}
\[
 w_{\mathrm{lead}}(x)=a\cos(\xi\cdot x+\varphi),\qquad a\cdot\xi=0.
\]
Here \NSi{NS-F-P011-024}{a\in\mathbb R^3} is a constant amplitude,
\NSi{NS-F-P011-025}{\xi\in\mathbb R^3} is the wavevector, and
\NSi{NS-F-P011-026}{\varphi\in\mathbb R} is a phase. The velocity oscillates orthogonally to its
wavevector. To separate distinct localized pulses while
\NSPage{012}%
retaining periodic averaging, we construct an extended field
\NSi{NS-F-P012-001}{\widetilde w(r,\theta,z,t,Y)} depending on an independent auxiliary variable
\NSi{NS-F-P012-002}{Y\in\mathbb T^2}. Its physical value is
\NSFormula{NS-F-P012-003}{display}{none}
\[
 w(r,\theta,z,t)=\widetilde w(r,\theta,z,t,Y(r,t)),
\]
where the fixed phase map \NSi{NS-F-P012-004}{Y(r,t)} is defined in \eqref{eq:6.3}. The covariance
is computed before this evaluation. We use angular and auxiliary averages
\NSFormula{NS-F-P012-005}{display}{none}
\[
 \langle F\rangle_\theta:=\frac1{2\pi}\int_0^{2\pi}F\,d\theta,
 \qquad
 \langle F\rangle_Y:=\int_{\mathbb T^2}F\,dY,
 \qquad
 \int_{\mathbb T^2}1\,dY=1.
\]
Their composition is denoted by
\NSi{NS-F-P012-006}{\langle F\rangle:=\langle\langle F\rangle_\theta\rangle_Y}. The averaging
conventions are collected in Equations~\eqref{eq:3.7} to~\eqref{eq:3.9}. Components are taken in
the cylindrical basis. Localized pulses with distinct labels whose slow space-time supports overlap
receive disjoint supports in \NSi{NS-F-P012-007}{Y}. Products between such pulses vanish pointwise,
including after evaluation at \NSi{NS-F-P012-008}{Y(r,t)} (Lemma~\ref{lem:6.1}). Harmonics and
corrections with the same label retain their interactions.

In the actual construction the amplitudes, phases, and cutoffs vary. We evaluate localized vector
potentials at \NSi{NS-F-P012-009}{Y(r,t)} to obtain
\NSi{NS-F-P012-010}{\mathcal A_{\mathrm{wave}}}, then take the physical curl:
\NSFormula{NS-F-P012-011}{display}{none}
\[
 w=\nabla\times\mathcal A_{\mathrm{wave}},\qquad \nabla\cdot w=0.
\]
All chain-rule terms from \NSi{NS-F-P012-012}{Y(r,t)} are included. Differentiating the amplitudes
and cutoffs produces smaller velocity terms, retained in Lemma~\ref{lem:7.7}.

The chosen phases have nonzero integer angular frequencies, so each wave has zero angular mean and
the cosine square has angular average \NSi{NS-F-P012-013}{1/2}. The two families produce the
momentum-flux directions used to choose \NSi{NS-F-P012-014}{c_1,c_2>0} above. These coefficients
become squared amplitude weights. The covariance construction, with its physical scaling and curl
corrections, gives
\NSFormula{NS-F-P012-015}{display}{none}
\[
 \begin{pmatrix}
  \langle\widetilde w_r\widetilde w_\theta\rangle\\
  \langle\widetilde w_r\widetilde w_z\rangle
 \end{pmatrix}
 =T+\text{higher-order terms}.
\]
The higher-order terms include the curl corrections and are smaller by positive powers of
\NSi{NS-F-P012-016}{q}, with derivative bounds proved in Corollary~\ref{cor:7.8}. The radial
divergence of the displayed leading covariance cancels the leading negative stress divergence of
the background (Propositions~\ref{prop:7.5} and~\ref{prop:9.5}). The extended residual still
contains oscillatory terms and angular means that depend on \NSi{NS-F-P012-017}{Y}; these are
evaluated at \NSi{NS-F-P012-018}{Y(r,t)} in the physical equation.

The same quadratic products govern energy transfer from the background. We consider the
homogeneous linearized model
\NSFormula{NS-F-P012-019}{display}{none}
\[
 \mathcal L_{u_B}(w_{\mathrm{lin}},\pi_{\mathrm{lin}})=0,
 \qquad \nabla\cdot w_{\mathrm{lin}}=0.
\]
Under sufficient spatial decay, integration by parts gives
\NSFormula{NS-F-P012-020}{display}{none}
\[
 \frac d{dt}\frac12\int|w_{\mathrm{lin}}|^2
 =-\sum_{i,j}\int(w_{\mathrm{lin}})_i(w_{\mathrm{lin}})_j\partial_j(u_B)_i
   -\int|\nabla w_{\mathrm{lin}}|^2.
\]
The first term transfers energy between the background and the perturbation; the second dissipates
it. The selected transverse velocity amplitudes extract energy from the shear and provide the two
momentum-flux directions whose positive span contains \NSi{NS-F-P012-021}{T}. During each pulse,
shear increases the magnitude of the radial component of the wavevector. Viscous damping
strengthens and eventually exceeds the amplification, so the pulse grows and then decays. The
temporal cutoff acts only in its exponentially small tails. The phase and amplitude equations
establish these properties in Lemmas~\ref{lem:7.1} and~\ref{lem:7.4}; Figure~\ref{fig:3}(b)
illustrates the resulting envelope.
